\documentclass[12pt]{article}

\usepackage{amsmath}
\usepackage[margin=1in]{geometry}
\usepackage{fancyhdr}
\usepackage{amsthm}
\newtheorem{lemma}{Lemma}
\usepackage{braket}
\usepackage{amssymb}
\usepackage{tikz}
\usepackage{bbm}
\usepackage{bm}
\usetikzlibrary{shapes.geometric}
\usepackage{enumerate}
\usepackage[shortlabels]{enumitem}

\newcommand{\bmalpha}{\bm{\alpha}}
\newcommand{\bmbeta}{\bm{\beta}}
\newcommand{\magnitude}[1]{\left\lvert #1 \right\rvert}

\pagestyle{fancy}
\fancyhead[l]{Hudson Benites}
\fancyhead[c]{Day 4}
\fancyfoot[c]{\thepage}
\renewcommand{\headrulewidth}{0.2pt}
\setlength{\headheight}{15pt}
\fancyhead[r]{\today}

\begin{document}
\begin{enumerate}[label=\Alph*., leftmargin=*]
    \item Let $\bmalpha:I\rightarrow \mathbb{R}^3$ be given by $\bmalpha(u)=\begin{cases}
        (u,u^2,0) \text{ if }x>0\\
        (u,0,u^2) \text{ if }x \leq 0
    \end{cases}.$ First, we show that $\bmalpha$ is always regular. When $x>0$, $\dot{\bmalpha}=(1,2u,0)$. When $x\leq 0$, $\dot{\bmalpha}=(1,0,2u)$. Therefore $\bmalpha$ is regular on its entire domain.
    Note that $\bm{T}=\frac{\dot{\bmalpha}}{\magnitude{\dot{\bmalpha}}}$. When $x>0$, we compute
    \begin{equation*}
        \bm{T}=\frac{1}{{(4u^2+1)}^{1/2}}\begin{bmatrix}
            1\\
            2u\\
            0
        \end{bmatrix}.
    \end{equation*}
    Note that $\bm{N}=\frac{\bm{T'}}{\magnitude{\bm{T'}}}$ also compute that\begin{equation*}
        \bm{N}=\frac{1}{{(4u^2+1)}^{1/2}}\begin{bmatrix}
            -2u\\
            1\\
            0
        \end{bmatrix}.
    \end{equation*} We compute $\bm{B}=\bm{T}\times\bm{N}$ as\begin{equation*}
        \bm{B}=(0+0,0+0, \frac{1}{4u^2+1}+\frac{4u^2}{4u^2+1})=(0,0,1).
    \end{equation*} By observation, $\frac{d\bm{B}}{ds}=0$. $\bm{N}$ is non-zero when $x>0$, and therefore $\tau = 0$.
    When $x\leq 0$, we compute\begin{equation*}
        \bm{T}=\frac{1}{{(4u^2+1)}^{1/2}}\begin{bmatrix}
            1\\
            0\\
            2u
        \end{bmatrix}.
    \end{equation*} We also compute that\begin{equation*}
        \bm{N}=\frac{1}{{(4u^2+1)}^{1/2}}\begin{bmatrix}
            -2u\\
            0\\
            1
        \end{bmatrix}.
    \end{equation*} Next, we compute $\bm{B}$.\begin{equation*}
        \bm{B}=(0-0, \frac{-4u^2-1}{4u^2+1}, 0-0)=(0, -1, 0).
    \end{equation*} Again, by observation, $\frac{d\bm{B}}{ds}=0$, and $\bm{N}$ is non-zero when $x\leq 0$, so $\tau = 0$. We can see that when $x>0$, $\bm{\alpha}$ lies completely in the $xy$-plane as it has no $z$ component, and when $x\leq0$, $\bm{\alpha}$ lies completely in the $xz$-plane, as it has no $y$ component.
\item By assumption, for every point $\bm{\alpha}(s)$, we have\begin{equation*}
    \bm{p}=\bm{\alpha}(s)+v(s)\bm{N}(s).
\end{equation*} Differentiating, we obtain\begin{equation*}
    0=\dot{\bm{\alpha}}+\dot{v}\bm{N}+\dot{\bm{N}}v.
\end{equation*}
Because $\bm{\alpha}$ is arc-length parameterized, $\dot{\bm{\alpha}}=\bm{T}$. The Serret-Frenet equations also tell us that $\frac{d\bm{N}}{ds}=-\kappa \bm{T} -\tau \bm{B}$. Substituting, we get that\begin{equation*}
    0=\bm{T}+\dot{v}\bm{N}+v(-\kappa \bm{T}-\tau \bm{B}).
\end{equation*} Combining like terms, we get\begin{equation*}
    0=(1-v\kappa)\bm{T}+\dot{v}\bm{N}-v\tau \bm{B}.
\end{equation*} After much struggle, we note that $\bm{T}, \bm{N} \text{and} \bm{B}$ are orthogonal, and thus the only the only way a linear
combination of them sum to 0 is if each of their coefficients are 0. Therefore, $\dot{v}=0$, and $v$ is constant. Also, $1-v\kappa=0\implies \kappa=\frac{1}{v}$. By assumption, $\kappa\neq 0$ everywhere,
and therefore $\tau=0$. Indeed, we have shown the curvature to be constant, the torsion to be 0. Next, consider $\magnitude{p-\bm{\alpha}}=\magnitude{vN}$. We know that $\magnitude{N}=1$, so $\magnitude{p-\bm{\alpha}}=\magnitude{v}$, which is constant.
Now, we have shown that $\alpha$ is a constant distance away from $p$, making it a sphere, which has torsion 0 and constant curvature. Thus, we are done.
\end{enumerate}
\end{document}